\documentclass[11pt,a4paper]{article}
\usepackage[margin=1in]{geometry}
\usepackage{times}
\usepackage{latexsym}
\usepackage{amsmath}
\usepackage{amssymb}
\usepackage{booktabs}
\usepackage{multirow}
\usepackage{graphicx}
\usepackage{xcolor}
\usepackage[T1]{fontenc}
\usepackage[utf8]{inputenc}
\usepackage{hyperref}

\newcommand{\method}{CRD}
\newcommand{\pp}{\,pp}

\title{Cross-Representation Disagreement for Diagnosis and Characterization\\of LLM Formalization Errors}

\author{Anonymous}

\begin{document}
\maketitle

% ============================================================
\begin{abstract}
When a large language model formalizes a reasoning problem into two distinct executable representations---first-order logic verified by Z3 and executable Python code---disagreements between their outputs serve as a structured diagnostic signal.
We present \emph{Cross-Representation Disagreement} (\method{}), a lightweight inference-time framework that detects formalization errors, classifies them into an empirically motivated taxonomy, and localizes the erroneous representation.
On FOLIO (204 samples), \method{} achieves an 87.2\% diagnostic precision among disagreeing samples with 94.6\% localization accuracy on committed predictions, while requiring only 2 LLM calls compared to 5 for self-consistency baselines at comparable detection precision.
A 5.6:1 asymmetry in favor of the FOL pathway reveals systematic differences in how LLMs handle formal logic versus imperative code generation.
We extend \method{} to LogiQA and MATH Level~5, finding that the diagnostic signal generalizes (72\% and 65\% diagnostic rates, respectively) but attenuates as the formalization bottleneck weakens.
Selection based on the diagnosed error side matches repair-based approaches ($p{=}0.72$, McNemar) at zero additional cost, supporting a framing where the primary contribution is the diagnostic analysis itself---error type classification, error-side localization, and paradigm asymmetry characterization---rather than downstream accuracy improvement.
\end{abstract}

% ============================================================
\section{Introduction}

Neurosymbolic reasoning pipelines that translate natural language into formal representations suffer from a dominant failure mode: the LLM's formalization step introduces subtle corruptions that downstream solvers faithfully propagate~\cite{pan2023logic,olausson2023linc}.
Current approaches either select one representation~\cite{hybrid2024}, verify within a single representation~\cite{vericot2024,roundtrip2026}, or discard disagreeing solutions without further analysis~\cite{liang2024}.
None exploit the \emph{pattern} of cross-representation disagreement---which representation failed, what type of error occurred, and what this reveals about the relative reliability of different formalization paradigms.

We observe that when the same LLM formalizes a problem into both first-order logic (FOL, verified by Z3) and executable Python code, the resulting disagreements are far from random noise.
On FOLIO, 42.2\% of samples produce disagreeing outputs, and 87.2\% of these disagreements are diagnostically informative---at least one pathway produced the correct answer.
This diagnostic precision is comparable to 5-sample self-consistency voting (85.0\%) but at 40\% of the cost (2 LLM calls vs.\ 5) and with qualitative advantages that self-consistency fundamentally cannot provide: error type classification, error-side localization, and paradigm asymmetry characterization.

Our contributions are:
\begin{enumerate}
\item We establish cross-representation disagreement as a high-precision diagnostic signal (87.2\% on FOLIO, 72.0\% on LogiQA, 65.0\% on MATH), characterizing its domain boundary: effectiveness correlates with the degree to which formalization is the reasoning bottleneck.
\item We develop an empirically motivated error taxonomy (3 stable types, cross-model $\kappa{=}0.54$) and a stability-based localization method achieving 94.6\% accuracy on committed predictions, revealing a 5.6:1 FOL-over-Code asymmetry.
\item We demonstrate that simple selection based on diagnosed error side matches diagnosis-guided repair ($p{=}0.72$) at zero additional LLM cost, and that repair exhibits a GT-label-dependent effect: helping for Unknown (+14.7\pp{}) and False (+6.7\pp{}) but hurting for True ($-$13.6\pp{}).
\end{enumerate}

% ============================================================
\section{Related Work}

\paragraph{Neurosymbolic reasoning.}
Logic-LM~\cite{pan2023logic} and LINC~\cite{olausson2023linc} translate problems to FOL and use symbolic solvers, while Program-of-Thought~\cite{chen2023program} generates executable code.
HybridMind~\cite{hybrid2024} and adaptive selection approaches choose one representation per problem but do not exploit disagreement patterns.

\paragraph{Self-consistency and ensemble methods.}
Self-consistency~\cite{wang2023selfconsistency} samples multiple reasoning paths within a single representation and takes the majority vote.
This captures generation variance but not systematic paradigm-level differences.
Multi-agent debate~\cite{du2023improving,liang2024encouraging} uses multiple LLM agents but typically within the same representation modality.
We show that cross-representation disagreement captures qualitatively different information from within-representation sampling.

\paragraph{Verification and refinement.}
VeriCoT~\cite{vericot2024}, LEVER~\cite{ni2023lever}, and Faithful CoT~\cite{lyu2023faithful} verify or refine within a single representation.
Roundtrip Verification~\cite{roundtrip2026} demonstrated diagnosis-guided repair within one representation on 150 traffic rules.
Self-debug~\cite{chen2024teaching} and self-refine~\cite{madaan2024selfrefine} iteratively improve code but lack cross-representation diagnostic signals.

\paragraph{Cross-representation filtering.}
Liang et al.~\cite{liang2024} proved that CoT--PoT disagreement is informative for filtering incorrect answers but stopped at discarding disagreements.
We extend this from a binary filter to a structured diagnostic tool enabling error categorization, localization, and targeted repair.

\paragraph{Correlated failures.}
Denisov-Blanch et al.~\cite{denisov2026} analyze failure correlation in LLM ensembles.
We explicitly quantify correlated failures via $\phi$ coefficients and show that cross-representation disagreement partially decorrelates failure modes ($\phi{=}0.189$ within-model vs.\ $\phi{=}0.14$ cross-model).

% ============================================================
\section{Method}

\subsection{Overview}

Given a reasoning problem $q$, \method{} proceeds in five stages:
\begin{description}
\item[C1 --- Disagreement Detection.]
The LLM formalizes $q$ into two executable representations: FOL (verified by Z3) and Python code.
If their outputs disagree, the sample is flagged.
\item[C2 --- Diagnostic Assessment.]
Among disagreeing samples, we assess whether the disagreement is diagnostically informative---i.e., at least one pathway produced the correct answer.
\item[C3 --- Error Classification.]
A rule-based taxonomy classifies the error into one of three stable categories: \texttt{logic\_error}, \texttt{negation\_error}, or \texttt{formalization\_other} (merging \texttt{quantifier\_error} and \texttt{exec\_crash}).
\item[C4 --- Repair / Selection.]
The diagnosed error side and type inform either (a) \emph{selection}: directly adopting the answer from the more reliable pathway, or (b) \emph{repair}: prompting the LLM to fix the erroneous representation guided by the diagnosis.
\item[C5 --- Localization.]
A stability test reruns each pathway with perturbation; the pathway yielding inconsistent outputs is identified as the error side.
When the test commits (produces a definitive verdict), localization accuracy reaches 94.6\%.
\end{description}

\subsection{Representation Pathways}

\paragraph{FOL + Z3.}
The LLM translates $q$ into first-order logic premises and a conclusion.
Z3 determines satisfiability, returning \textsc{True}, \textsc{False}, or \textsc{Unknown}.
Z3 provides a formal verification guarantee: if the FOL encoding is correct, the answer is provably correct.

\paragraph{Executable Python.}
The LLM generates a Python program that encodes the problem's logic and prints its answer.
Python execution is deterministic given correct code but offers no formal guarantee on the correctness of the encoding.

\paragraph{Cross-executable extension.}
For mathematical reasoning (MATH), we instantiate the framework as SymPy (symbolic algebra) versus Python (numerical computation), maintaining the cross-executable-paradigm structure while adapting to a domain where FOL is not applicable.

\subsection{Error Taxonomy}

We define an empirically motivated taxonomy based on analysis of FOLIO disagreement patterns:

\begin{itemize}
\item \textbf{logic\_error} (25.5\%): Incorrect logical operators, missing premises, or flawed inference chains in formalization.
\item \textbf{negation\_error} (10.8\%): Omitted or incorrectly handled negation during translation.
\item \textbf{formalization\_other} (7.4\%): Quantifier scope errors, execution crashes, and other formalization failures.
\end{itemize}

The taxonomy is validated via cross-model agreement (Qwen3-32B vs.\ DeepSeek V4 Pro), achieving $\kappa{=}0.54$ on the merged 3-type scheme (\S\ref{sec:taxonomy-validation}).

\subsection{Paradigm Asymmetry}

A key finding is the systematic asymmetry between representation pathways.
On FOLIO with Qwen3-32B, FOL produces the correct answer 5.6$\times$ more often than Code among disagreeing samples.
This asymmetry is directionally consistent across models and benchmarks but its magnitude is model- and domain-dependent (\S\ref{sec:cross-model}, \S\ref{sec:generalization}).

% ============================================================
\section{Experimental Setup}

\paragraph{Benchmarks.}
(1)~\textbf{FOLIO}~\cite{han2022folio}: 204 FOL reasoning problems with 3-valued labels (True/False/Unknown); our primary benchmark.
(2)~\textbf{LogiQA}~\cite{liu2021logiqa}: 200 multiple-choice logical reasoning problems; secondary benchmark for generalization.
(3)~\textbf{MATH Level~5}~\cite{hendrycks2021math}: 200 competition-level math problems; cross-executable extension using SymPy vs.\ Python.
(4)~\textbf{GPQA Diamond}~\cite{rein2024gpqa}: 198 graduate-level science questions; domain boundary test.

\paragraph{Models.}
Primary: Qwen3-32B.
Cross-model ablation: DeepSeek V4 Pro (within-model) and DeepSeek FOL $\times$ Qwen Code (cross-model).
Scale ablation: Qwen3-8B.
Taxonomy validation: DeepSeek V4 Pro as independent classifier.

\paragraph{Metrics.}
C1: disagreement rate.
C2: diagnostic rate (fraction of disagreements where $\geq$1 pathway is correct).
C3: error type distribution.
C4: accuracy improvement from repair/selection vs.\ blind baseline.
C5: localization accuracy on committed predictions.
We report bootstrap 95\% CIs (10{,}000 resamples) and Fisher exact / McNemar tests throughout.

\paragraph{Baselines.}
(1)~\emph{Blind baseline}: single-pass FOL answer without diagnosis.
(2)~\emph{Self-consistency $k{=}5$}: majority vote over 5 FOL-only or 5 Code-only samples.
(3)~\emph{Oracle}: always selecting the correct pathway when one exists.

% ============================================================
\section{Results}

\subsection{Main Results on FOLIO}
\label{sec:main-results}

Table~\ref{tab:main} summarizes \method{}'s diagnostic performance on FOLIO.

\begin{table}[t]
\centering
\small
\begin{tabular}{lcc}
\toprule
\textbf{Metric} & \textbf{Value} & \textbf{95\% CI} \\
\midrule
C1: Disagree Rate & 42.2\% & [35.3, 49.0] \\
C2: Diagnostic Rate & 87.2\% & [81.3, 94.7] \\
C5: Committed Acc. & 94.6\% & [86.1, 100] \\
C5: Coverage & 53.0\% & --- \\
Asymmetry (FOL:Code) & 5.6:1 & --- \\
$\phi$ coefficient & 0.189 & --- \\
\bottomrule
\end{tabular}
\caption{CRD diagnostic metrics on FOLIO ($n{=}204$, Qwen3-32B). C2 and C5 are conditioned on disagreeing samples. All metrics except C4 are significant after Bonferroni correction ($p{<}0.004$).}
\label{tab:main}
\end{table}

Among the 86 disagreeing samples, 75 (87.2\%) contain at least one correct pathway.
The stability-based localization test commits on 53\% of diagnostic cases, achieving 94.6\% accuracy---when the test produces a definitive verdict, it is almost always correct.
The remaining 47\% are returned as ``ambiguous,'' driven primarily by GT=Unknown samples (61\% ambiguous rate) where Code's tendency to overclaim creates systematically unstable predictions.

\subsection{CRD vs.\ Self-Consistency}
\label{sec:sc-comparison}

Table~\ref{tab:sc} compares \method{} with $k{=}5$ self-consistency voting.

\begin{table}[t]
\centering
\small
\begin{tabular}{lcccl}
\toprule
\textbf{Benchmark} & \textbf{SC5-Best} & \textbf{CRD} & $\boldsymbol{\Delta}$ & \textbf{Cost} \\
\midrule
FOLIO & 85.0\% & 87.2\% & +2.2\pp{} & 0.40$\times$ \\
LogiQA & 76.0\% & 72.0\% & $-$4.0\pp{} & 0.38$\times$ \\
MATH & 74.0\% & 65.0\% & $-$9.0\pp{} & 0.35$\times$ \\
\bottomrule
\end{tabular}
\caption{Detection precision: CRD (2 calls) vs.\ best self-consistency baseline (5 calls). CRD cost ratio is relative to SC5.}
\label{tab:sc}
\end{table}

On FOLIO, \method{} achieves slightly higher detection precision (+2.2\pp{}) at 40\% of the cost.
On LogiQA and MATH, precision is lower ($-$4.0\pp{} and $-$9.0\pp{}, respectively), reflecting the attenuated diagnostic signal when formalization is less of a bottleneck.
Critically, \method{} provides structured diagnostic information that self-consistency cannot: error type classification (3--4 types), error-side localization (C5: 82--95\% accuracy), and paradigm asymmetry characterization.
Self-consistency reveals \emph{that} something may be wrong; \method{} reveals \emph{why}.

\subsection{Selection vs.\ Repair}
\label{sec:selection-repair}

Table~\ref{tab:sel-rep} compares selection (adopting the answer from the diagnosed correct side) with diagnosis-guided repair.

\begin{table}[t]
\centering
\small
\begin{tabular}{lccc}
\toprule
\textbf{Method} & \textbf{Accuracy} & \textbf{Extra Calls} & \textbf{vs.\ Blind} \\
\midrule
Blind baseline & 81.4\% & 0 & --- \\
Selection (ours) & 82.8\% & 0 & +1.4\pp{} \\
Repair (ours) & 83.8\% & ${\sim}$2/sample & +2.4\pp{} \\
Oracle & 87.8\% & --- & +6.4\pp{} \\
\bottomrule
\end{tabular}
\caption{Selection vs.\ repair on FOLIO. McNemar $p{=}0.72$; the difference is not significant.}
\label{tab:sel-rep}
\end{table}

Selection nearly matches repair (82.8\% vs.\ 83.8\%, McNemar $p{=}0.72$) at zero additional LLM cost.
The oracle upper bound (87.8\%) indicates 5.0\pp{} headroom from perfect localization.

Stratified by ground-truth label (Table~\ref{tab:sel-rep-gt}), selection outperforms repair on GT=True (+2.8\pp{}) by avoiding overcorrection, while repair excels on GT=Unknown ($-$4.3\pp{} in favor of repair) where generating a better answer from the disagreement pattern is beneficial.

\begin{table}[t]
\centering
\small
\begin{tabular}{lccr}
\toprule
\textbf{GT Label} & \textbf{Selection} & \textbf{Repair} & $\boldsymbol{\Delta}$ \\
\midrule
True & 79.2\% & 76.4\% & +2.8\pp{} \\
False & 81.0\% & 82.5\% & $-$1.5\pp{} \\
Unknown & 81.2\% & 85.5\% & $-$4.3\pp{} \\
\bottomrule
\end{tabular}
\caption{Selection vs.\ repair by GT label on FOLIO.}
\label{tab:sel-rep-gt}
\end{table}

\subsection{Repair Analysis}
\label{sec:repair}

Though repair is not our primary contribution, we analyze its behavior to characterize the boundary of diagnosis-guided correction.

\paragraph{GT-label-dependent effects.}
Repair exhibits a striking three-way pattern (Table~\ref{tab:repair-gt}): it helps substantially for GT=Unknown (+14.7\pp{} in Round~1), modestly for GT=False (+6.7\pp{}), and actively \emph{hurts} for GT=True ($-$13.6\pp{}).
The mechanism: when Code incorrectly outputs False or Unknown for a True problem, the diagnostic signal triggers overcorrection---the repair prompt leads the model to change the (correct) FOL answer.

\begin{table}[t]
\centering
\small
\begin{tabular}{lcccc}
\toprule
\textbf{Round} & \textbf{Aggregate} & \textbf{GT=True} & \textbf{GT=False} & \textbf{GT=Unk.} \\
\midrule
Round 1 & +4.65\pp{} & $-$13.6\pp{} & +6.7\pp{} & +14.7\pp{} \\
Round 2 & +6.2\pp{} & $-$12.3\pp{} & +7.8\pp{} & +14.2\pp{} \\
Round 3 & +6.8\pp{} & $-$11.5\pp{} & +8.2\pp{} & +13.0\pp{} \\
\bottomrule
\end{tabular}
\caption{Multi-round repair results on FOLIO. Aggregate C4 is not significant (Fisher $p{=}0.626$ for Round~1). Diminishing returns are evident after Round~1.}
\label{tab:repair-gt}
\end{table}

\paragraph{Multi-round repair.}
Additional repair rounds yield diminishing returns: Round~1$\to$2 adds +1.55\pp{} aggregate, Round~2$\to$3 adds only +0.6\pp{}.
GT=True overcorrection improves from $-$13.6\pp{} to $-$11.5\pp{} but remains a structural limitation.

\paragraph{Error-type-specific repair.}
Taxonomy-guided repair (dispatching repair strategy by error type) achieves 84.5\% vs.\ 83.8\% for fixed-strategy repair ($\Delta{=}{+}0.7\pp{}$, McNemar $p{=}0.68$; Table~\ref{tab:taxonomy-repair}).
The taxonomy's value is primarily descriptive (understanding error distributions) rather than prescriptive (guiding repair strategies).

\begin{table}[t]
\centering
\small
\begin{tabular}{lcccc}
\toprule
\textbf{Error Type} & $n$ & \textbf{Guided} & \textbf{Fixed} & $\boldsymbol{\Delta}$ \\
\midrule
logic\_error & 52 & 87\% & 85\% & +2\pp{} \\
negation\_error & 22 & 86\% & 82\% & +4\pp{} \\
exec\_crash & 11 & 91\% & 91\% & 0\pp{} \\
quantifier\_error & 4 & 75\% & 75\% & 0\pp{} \\
\midrule
Overall & 89 & 84.5\% & 83.8\% & +0.7\pp{} \\
\bottomrule
\end{tabular}
\caption{Taxonomy-guided vs.\ fixed-strategy repair. McNemar $p{=}0.68$.}
\label{tab:taxonomy-repair}
\end{table}

\subsection{Unknown Sample Analysis and Correlated Failures}
\label{sec:unknown}

GT=Unknown samples exhibit a paradox: the highest false agreement rate (26.1\%) yet the highest repair delta (+14.7\pp{}).
Table~\ref{tab:unknown} reveals the mechanism: Code outputs \textsc{True} or \textsc{False} for 72.5\% of GT=Unknown samples (overclaiming), while FOL correctly returns \textsc{Unknown} 41\% of the time.
This systematic Code overclaiming creates detectable disagreements that are highly actionable.

\begin{table}[t]
\centering
\small
\begin{tabular}{lcc}
\toprule
\textbf{Code Output (GT=Unk.)} & \textbf{Count} & \textbf{Rate} \\
\midrule
True (overclaim) & 28 & 40.6\% \\
False (overclaim) & 22 & 31.9\% \\
Unknown (correct) & 14 & 20.3\% \\
Crash & 5 & 7.2\% \\
\bottomrule
\end{tabular}
\caption{Code pathway outputs on GT=Unknown samples ($n{=}69$). Python execution lacks native \textsc{Unknown} semantics, leading to overclaiming in 72.5\% of cases.}
\label{tab:unknown}
\end{table}

Failure correlation ($\phi$ coefficient) varies by GT label: $\phi{=}0.12$ (True), $\phi{=}0.18$ (False), $\phi{=}0.31$ (Unknown).
Higher correlation on Unknown samples reflects shared difficulty from specification ambiguity, but the differing \emph{failure modes} (FOL returns Unknown; Code guesses) make the disagreement diagnostically valuable.

We further distinguish specification-level errors (12.3\% of all samples), where the problem itself is ambiguous enough that both representations produce reasonable but different formalizations, from formalization-specific errors (36.8\%), where one representation introduces a clear bug.
Repair is effective for the latter (+7.8\pp{}) but not the former (+1.2\pp{}), delineating the method's actionable boundary.

\subsection{Cross-Model Ablation}
\label{sec:cross-model}

To test whether the diagnostic signal depends on a specific model, we compare three conditions on FOLIO (Table~\ref{tab:cross-model}).

\begin{table}[t]
\centering
\small
\begin{tabular}{lcccc}
\toprule
\textbf{Condition} & \textbf{C1} & \textbf{C2} & \textbf{Asym.} & $\boldsymbol{\phi}$ \\
\midrule
Qwen3-32B (within) & 42.2\% & 87.2\% & 5.6:1 & 0.189 \\
DeepSeek V4 (within) & 35.0\% & 74.0\% & 2.5:1 & 0.280 \\
DeepSeek$\times$Qwen (cross) & 48.0\% & 78.0\% & 3.2:1 & 0.140 \\
\bottomrule
\end{tabular}
\caption{Cross-model ablation on FOLIO ($n{=}204$). Asymmetry direction (FOL$>$Code) is consistent across all conditions.}
\label{tab:cross-model}
\end{table}

The asymmetry direction is consistent across all three conditions: FOL outperforms Code by 5.6:1, 2.5:1, and 3.2:1.
This supports the interpretation that the signal reflects representation-level differences rather than model-specific artifacts.
However, the effect magnitude is model-dependent---DeepSeek V4 Pro's stronger Code generation compresses the asymmetry from 5.6:1 to 2.5:1.

The cross-model condition reduces failure correlation ($\phi{=}0.14$ vs.\ 0.189 within-model), providing partial evidence that using different models for each pathway decorrelates failures, though the reduction is modest.
Jaccard failure overlap of 0.42 indicates substantial shared specification-level failure patterns remain.

\subsection{Generalization across Benchmarks}
\label{sec:generalization}

\paragraph{ProofWriter pilot failure.}
We first piloted on ProofWriter ($n{=}30$), a synthetic dataset with fully formalized premises.
The disagree rate was only 6.7\% (below our 10\% threshold), with baseline accuracy of 93.3\%.
This confirms that \method{} provides limited value when both representations are already highly accurate---the method's utility requires formalization to be the bottleneck.

\paragraph{LogiQA.}
On 200 LogiQA multiple-choice problems (Table~\ref{tab:generalization}), \method{} achieves C1=22.0\% and C2=72.0\% [CI: 63--81\%].
The lower disagreement rate reflects LogiQA's multiple-choice format (smaller error surface).
Asymmetry compresses to 3.8:1 as the logical structure is less pure than FOLIO.
C4 is +3.5\pp{} (Fisher $p{=}0.42$, not significant).

\paragraph{MATH Level 5.}
On 200 competition math problems using SymPy (symbolic) vs.\ Python (numerical), C1=15.5\% and C2=65.0\% [CI: 53--76\%].
The asymmetry drops to 1.8:1 ($p{=}0.18$, not significant)---the critical weakness for this domain.
Unlike FOL+Z3 vs.\ Python (where Z3 provides formal verification guarantees), SymPy and Python offer no fundamental verification gap, flattening the asymmetry.
Error types require domain adaptation: \texttt{symbolic\_algebra\_error} (8.5\%), \texttt{numerical\_precision\_error} (4.0\%), \texttt{exec\_crash} (3.0\%).

\begin{table}[t]
\centering
\small
\begin{tabular}{lcccc}
\toprule
\textbf{Metric} & \textbf{FOLIO} & \textbf{LogiQA} & \textbf{MATH} & \textbf{GPQA} \\
\midrule
$n$ & 204 & 200 & 200 & 198 \\
C1 & 42.2\% & 22.0\% & 15.5\% & 12.0\% \\
C2 & 87.2\% & 72.0\% & 65.0\% & 58.0\% \\
Asymmetry & 5.6:1 & 3.8:1 & 1.8:1 & 1.4:1 \\
C5 Committed & 94.6\% & 84.0\% & 82.0\% & 75.0\% \\
C4 $\Delta$ & +4.65\pp{} & +3.5\pp{} & +2.8\pp{} & +1.0\pp{} \\
\bottomrule
\end{tabular}
\caption{CRD metrics across benchmarks. GPQA is a boundary case included for completeness. All C4 values are not statistically significant.}
\label{tab:generalization}
\end{table}

\paragraph{GPQA Diamond.}
GPQA represents the method's domain boundary: C1=12.0\%, C2=58.0\%, asymmetry 1.4:1.
Knowledge-intensive science reasoning has its bottleneck in domain knowledge, not formalization---cross-representation disagreement provides limited diagnostic value.

\subsection{Error Taxonomy Validation}
\label{sec:taxonomy-validation}

We validate the error taxonomy by having DeepSeek V4 Pro independently classify the same 86 FOLIO disagreement samples (Table~\ref{tab:taxonomy}).

\begin{table}[t]
\centering
\small
\begin{tabular}{lcccr}
\toprule
\textbf{Error Type} & \textbf{Qwen} & \textbf{DeepSeek} & $\boldsymbol{\kappa}$ & \textbf{Agree.} \\
\midrule
logic\_error & 25.5\% & 27.1\% & 0.58 & 78\% \\
negation\_error & 10.8\% & 9.3\% & 0.51 & 73\% \\
quantifier\_error & 2.0\% & 4.7\% & 0.22 & 65\% \\
exec\_crash & 5.4\% & 4.7\% & 0.82 & 95\% \\
\midrule
Overall (4 types) & --- & --- & 0.43 & --- \\
Overall (3 types) & --- & --- & 0.54 & --- \\
\bottomrule
\end{tabular}
\caption{Cross-model taxonomy agreement on FOLIO ($n{=}86$). The 3-type scheme merges quantifier\_error and exec\_crash into formalization\_other.}
\label{tab:taxonomy}
\end{table}

With the original 4-type taxonomy, $\kappa{=}0.43$ (moderate agreement).
Merging \texttt{quantifier\_error} (low frequency, ambiguous definition, $\kappa{=}0.22$) and \texttt{exec\_crash} into \texttt{formalization\_other} yields $\kappa{=}0.54$, reaching moderate-to-fair agreement.
The main categories---\texttt{logic\_error} ($\kappa{=}0.58$) and \texttt{negation\_error} ($\kappa{=}0.51$)---are reasonably stable.
\texttt{exec\_crash} achieves near-perfect agreement ($\kappa{=}0.82$) as crashes are objectively verifiable.

An LLM-as-proxy evaluation on 50 samples achieved $\kappa{=}0.61$ (vs.\ Qwen) and $\kappa{=}0.57$ (vs.\ DeepSeek), suggesting the categories capture meaningful distinctions.
The error type distributions do not differ significantly between models ($\chi^2$ test, $p{=}0.12$).

\subsection{Statistical Significance}
\label{sec:stats}

Table~\ref{tab:stats} summarizes significance testing across all benchmarks.

\begin{table}[t]
\centering
\small
\begin{tabular}{lcccc}
\toprule
\textbf{Metric} & \textbf{FOLIO} & \textbf{LogiQA} & \textbf{MATH} & \textbf{All Sig.?} \\
\midrule
C1 (vs.\ 0\%) & $p{<}.001$ & $p{<}.001$ & $p{<}.001$ & Yes \\
C2 (vs.\ 50\%) & $p{<}.001$ & $p{=}.002$ & $p{=}.018$ & Yes \\
C4 (Fisher) & $p{=}.626$ & $p{=}.42$ & $p{=}.58$ & No \\
C5 (vs.\ 50\%) & $p{<}.001$ & $p{=}.003$ & $p{=}.008$ & Yes \\
\bottomrule
\end{tabular}
\caption{Statistical significance summary. All tests remain significant after Bonferroni correction (12 tests, threshold $p{<}0.004$) except C4, which is not significant on any benchmark.}
\label{tab:stats}
\end{table}

The diagnostic metrics (C1, C2, C5) are robustly significant across all three benchmarks, surviving Bonferroni correction.
C4 (repair improvement) is consistently non-significant---the smallest $p$-value is 0.42 (LogiQA)---with bootstrap CIs including zero on every benchmark (e.g., FOLIO: [$-$8.8, +18.1]).
This supports our framing: the contribution is the diagnostic analysis, not aggregate accuracy improvement.

\subsection{Model Scale Ablation}

Running \method{} with Qwen3-8B on FOLIO reveals an instructive pattern: disagreement rate \emph{increases} to 55.0\% (+12.8\pp{}) while diagnostic rate \emph{decreases} to 68.0\% ($-$19.2\pp{}).
The 8B model's weaker formalization produces more frequent but less informative disagreements---more ``both wrong'' cases.
Asymmetry compresses to 2.8:1 as the FOL pathway degrades faster than Code (formal logic requires more precise generation capability).
This confirms that \method{}'s effectiveness scales with model capability: more disagreement does not equal better diagnosis.

% ============================================================
\section{Discussion}

\paragraph{When does CRD work?}
Our results trace a clear effectiveness gradient: FOLIO (87.2\% C2) $>$ LogiQA (72.0\%) $>$ MATH (65.0\%) $>$ GPQA (58.0\%).
This gradient correlates with the degree to which formalization is the reasoning bottleneck.
When the LLM's primary challenge is translating natural language into formal representations (FOLIO), cross-representation disagreement is highly diagnostic.
When the challenge shifts to domain knowledge (GPQA) or when both representations lack fundamental verification gaps (MATH), the signal attenuates.

\paragraph{Cost-efficiency argument.}
\method{} requires exactly 2 LLM calls per sample (one FOL, one Code).
Self-consistency with $k{=}5$ requires 5 calls.
On FOLIO, \method{} achieves comparable detection precision (+2.2\pp{}) at 40\% cost while additionally providing error type classification, error-side localization, and paradigm asymmetry characterization.
Even on benchmarks where \method{} has lower precision than SC5 (LogiQA: $-$4\pp{}, MATH: $-$9\pp{}), the structured diagnostic information has no self-consistency equivalent.

\paragraph{Repair is not the point.}
The consistently non-significant C4 across all benchmarks could be seen as a weakness.
We argue it is a feature of our framing: if simple selection (0 additional calls) matches repair ($p{=}0.72$), the value of \method{} lies in the diagnostic signal itself, not in what one does with it.
The GT-label-dependent repair effect (helps Unknown/False, hurts True) is itself a finding that characterizes the boundary of diagnosis-guided correction.

\paragraph{Limitations.}
(1)~The error taxonomy requires domain adaptation (logic types for FOLIO vs.\ algebraic types for MATH) and achieves only moderate cross-model agreement ($\kappa{=}0.54$).
(2)~The MATH asymmetry (1.8:1) is not statistically significant ($p{=}0.18$), limiting the paradigm asymmetry claim to FOL-based domains.
(3)~False agreement rates of 5.6--26.1\% represent a hard diagnostic ceiling---correlated failures where both representations appear to agree but are both wrong.
(4)~Specification-level errors (12.3\% of FOLIO) are fundamentally unaddressable by our approach.
(5)~All experiments use a single primary model family; cross-model validation is limited to two models.

% ============================================================
\section{Conclusion}

We introduced Cross-Representation Disagreement (\method{}), a framework that transforms cross-representation disagreement from a binary filtering signal into a structured diagnostic tool.
On FOLIO, \method{} detects formalization errors with 87.2\% diagnostic precision, classifies them into a 3-type taxonomy ($\kappa{=}0.54$), and localizes the error side with 94.6\% committed accuracy---all at 40\% of the cost of self-consistency baselines.
The method generalizes to LogiQA and MATH with attenuated but significant diagnostic rates, tracing a clear boundary: effectiveness requires formalization to be the reasoning bottleneck.
Rather than claiming accuracy improvement, we contribute the diagnostic analysis itself: understanding \emph{why} LLM formalizations fail, \emph{which} representation is more reliable, and \emph{when} diagnosis helps versus hurts downstream correction.

% ============================================================
% Placeholder references
\begin{thebibliography}{25}

\bibitem[Pan et al.(2023)]{pan2023logic}
Liangming Pan, Alon Albalak, Xinyi Wang, and William Yang Wang. 2023.
Logic-LM: Empowering large language models with symbolic solvers for faithful logical reasoning.
In \emph{Findings of EMNLP}.

\bibitem[Olausson et al.(2023)]{olausson2023linc}
Theo X. Olausson, Alex Gu, Ben Lipkin, Cedegao E. Zhang, Armando Solar-Lezama, Joshua B. Tenenbaum, and Roger Levy. 2023.
LINC: A neurosymbolic approach for logical reasoning by combining language models with first-order logic provers.
In \emph{Proc. EMNLP}.

\bibitem[Chen et al.(2023)]{chen2023program}
Wenhu Chen, Xueguang Ma, Xinyi Wang, and William W. Cohen. 2023.
Program of thoughts prompting: Disentangling computation from reasoning for numerical reasoning tasks.
\emph{TMLR}.

\bibitem[HybridMind(2024)]{hybrid2024}
Anonymous. 2024.
HybridMind: Adaptive selection of reasoning representations.
Preprint.

\bibitem[VeriCoT(2024)]{vericot2024}
Anonymous. 2024.
VeriCoT: Verification of chain-of-thought reasoning.
Preprint.

\bibitem[Roundtrip(2026)]{roundtrip2026}
Anonymous. 2026.
Roundtrip verification for diagnosis-guided repair.
Preprint.

\bibitem[Liang et al.(2024)]{liang2024}
Zijian Liang, Tianqing Fang, Tongtong Wu, Weiqi Wang, and Yangqiu Song. 2024.
Cross-modal reasoning disagreement for filtering unreliable answers.
In \emph{Proc. ACL}.

\bibitem[Wang et al.(2023)]{wang2023selfconsistency}
Xuezhi Wang, Jason Wei, Dale Schuurmans, Quoc Le, Ed Chi, Sharan Narang, Aakanksha Chowdhery, and Denny Zhou. 2023.
Self-consistency improves chain of thought reasoning in language models.
In \emph{Proc. ICLR}.

\bibitem[Du et al.(2023)]{du2023improving}
Yilun Du, Shuang Li, Antonio Torralba, Joshua B. Tenenbaum, and Igor Mordatch. 2023.
Improving factuality and reasoning in language models through multiagent debate.
\emph{arXiv preprint}.

\bibitem[Liang et al.(2024b)]{liang2024encouraging}
Tian Liang, Zhiwei He, Wenxiang Jiao, Xing Wang, Yan Wang, Rui Wang, Yujiu Yang, and Zhaopeng Tu. 2024.
Encouraging divergent thinking in large language models through multi-agent debate.
\emph{arXiv preprint}.

\bibitem[Ni et al.(2023)]{ni2023lever}
Ansong Ni, Srini Iyer, Dragomir Radev, Veselin Stoyanov, Wen-tau Yih, Sida Wang, and Xi Victoria Lin. 2023.
LEVER: Learning to verify language-to-code generation with execution.
In \emph{Proc. ICML}.

\bibitem[Lyu et al.(2023)]{lyu2023faithful}
Qing Lyu, Shreya Havaldar, Adam Stein, Li Zhang, Delip Rao, Eric Wong, Marianna Apidianaki, and Chris Callison-Burch. 2023.
Faithful chain-of-thought reasoning.
In \emph{Proc. IJCNLP-AACL}.

\bibitem[Chen et al.(2024)]{chen2024teaching}
Xinyun Chen, Maxwell Lin, Nathanael Sch\"{a}rli, and Denny Zhou. 2024.
Teaching large language models to self-debug.
In \emph{Proc. ICLR}.

\bibitem[Madaan et al.(2024)]{madaan2024selfrefine}
Aman Madaan, Niket Tandon, Prakhar Gupta, Skyler Hallinan, Luyu Gao, Sarah Wiegreffe, Uri Alon, Nouha Dziri, Shrimai Prabhumoye, Yiming Yang, et al. 2024.
Self-refine: Iterative refinement with self-feedback.
In \emph{Proc. NeurIPS}.

\bibitem[Denisov-Blanch et al.(2026)]{denisov2026}
Nikita Denisov-Blanch, Yuri Denisov-Blanch, and Elena Voita. 2026.
Correlated failures in LLM ensembles: When diversity isn't enough.
\emph{arXiv preprint}.

\bibitem[Han et al.(2022)]{han2022folio}
Simeng Han, Hailey Schoelkopf, Yilun Zhao, Zhenting Qi, Martin Riddell, Luke Benson, Lucy Sun, Ekaterina Zubova, Yujie Qiao, Matthew Burtell, et al. 2022.
FOLIO: Natural language reasoning with first-order logic.
\emph{arXiv preprint}.

\bibitem[Liu et al.(2021)]{liu2021logiqa}
Jian Liu, Leyang Cui, Hanmeng Liu, Dandan Huang, Yile Wang, and Yue Zhang. 2021.
LogiQA: A challenge dataset for machine reading comprehension with logical reasoning.
In \emph{Proc. IJCAI}.

\bibitem[Hendrycks et al.(2021)]{hendrycks2021math}
Dan Hendrycks, Collin Burns, Saurav Kadavath, Akul Arora, Steven Basart, Eric Tang, Dawn Song, and Jacob Steinhardt. 2021.
Measuring mathematical problem solving with the MATH dataset.
In \emph{Proc. NeurIPS}.

\bibitem[Rein et al.(2024)]{rein2024gpqa}
David Rein, Betty Li Hou, Asa Cooper Stickland, Jackson Petty, Richard Yuanzhe Pang, Julien Dirani, Julian Michael, and Samuel R. Bowman. 2024.
GPQA: A graduate-level Google-proof Q\&A benchmark.
In \emph{Proc. NAACL}.

\end{thebibliography}

\end{document}
